% Môn: Vật lý
% Lớp: 12
% Chương: Chương II. Khí lí tưởng
% Bài: L12C2 Bài 8. Mô hình động học phân tử chất khí · Đặc điểm chất rắn
% ID: L12C1B1-100-TN
% Mức: NB
% ID: L12C1B1-100-TN
% ID: L12C1B1-100-TN
% Mức: NB
% ID: L12C1B1-100-TN
% ID: L12C1B1-100-TN
% Mức: NB
% ID: L12C1B1-100-TN
% ID: L12C1B1-100-TN
% Mức: NB
% ID: L12C2B8-90-TN
%=== Câu 1 ===%
\dangbt{Đặc điểm chất rắn}
\begin{ex}
% ID: L12C2B8-90-TN
% Mức: NB
	{\color{red}\footnotesize\ttfamily [L12C1B1-TN-3]}\par %HienID
	Phát biểu nào sau đây là \textbf{sai} khi nói về chất rắn vô định hình?
	\choice
	{Vật rắn vô định hình không có cấu trúc tinh thể}
	{Vật rắn vô định hình không có nhiệt độ nóng chảy xác định}
	{Vật rắn vô định hình không có dạng hình học xác định}
	{\True Thuỷ tinh, nhựa đường, cao su, muối ăn là những chất rắn vô định hình}
	\loigiai{Thuỷ tinh, nhựa đường, cao su là những chất rắn vô định hình. Muối ăn là chất rắn kết tinh.\\
		Chất rắn vô định hình là chất rắn không có cấu trúc mạng tinh thể nên không có dạng hình học xác định như thủy tinh, nhựa, socola,...Muối ăn không phải chất rắn vô định hình sách giáo khoa có hình cấu trúc tinh thể NaCl}
\end{ex}

% ID: L12C1B1-101-TN
% Mức: NB
% ID: L12C1B1-101-TN
% ID: L12C1B1-101-TN
% Mức: NB
% ID: L12C1B1-101-TN
% ID: L12C1B1-101-TN
% Mức: NB
% ID: L12C1B1-101-TN
% ID: L12C1B1-101-TN
% Mức: NB
% ID: L12C2B8-91-TN
%=== Câu 2 ===%
\begin{ex}
% ID: L12C2B8-91-TN
% Mức: NB
	{\color{red}\footnotesize\ttfamily [L12C1B1-TN-16]}\par %HienID
	Phát biểu nào sau đây là \textbf{sai} khi nói về tính chất mạng tinh thể của chất rắn?
	\choice
	{Chất rắn kết tinh có cấu trúc mạng tinh thể xác định}
	{Cấu trúc mạng tinh thể khác nhau thì có tính chất của chất kết tinh khác nhau}
	{Các chất khác nhau có mạng tinh thể khác nhau}
	{\True Cùng một chất mạng tinh thể phải giống nhau} % Đáp án đúng
	\loigiai{
		Ta thấy cùng một chất cacbon (C) nhưng cấu trúc mạng tinh thể khác nhau thì chất cũng khác nhau như kim cương và than đá}
\end{ex}

% ID: L12C1B1-61-TN
% Mức: NB
% ID: L12C1B1-61-TN
% ID: L12C1B1-61-TN
% Mức: NB
% ID: L12C1B1-61-TN
% ID: L12C1B1-61-TN
% Mức: NB
% ID: L12C1B1-61-TN
% ID: L12C1B1-61-TN
% Mức: NB
% ID: L12C2B8-92-TN
%=== Câu 14 ===%
\begin{ex}
% ID: L12C2B8-92-TN
% Mức: NB
 Phát biểu nào sau đây đúng?
\choice
{Chất rắn vô định hình luôn có nhiệt độ nóng chảy xác định.}
{Nhiệt lượng và nhiệt độ là hai đại lượng hoàn toàn giống nhau.}
{Mọi quá trình nhiệt học đều không phụ thuộc điều kiện thực hiện.}
{\True Chất rắn kết tinh có cấu trúc tinh thể và nhiệt độ nóng chảy xác định.}
\loigiai{
Phát biểu đúng là: Chất rắn kết tinh có cấu trúc tinh thể và nhiệt độ nóng chảy xác định. Các phương án còn lại trái với mô hình, định nghĩa hoặc hệ thức nhiệt học tương ứng.
}
\end{ex}
